\section{A ranked equation language}

The Eq/$D$ grammar has no branching at the formula level: every formula
has a unique shape $D^k(s=t)$ for an integer $k\geq0$. Each $D$ binds two
variables, so the endpoints $s,t$ are terms in a context with $2k$ additional
bound variables. Store this formula as the triple $(k,s,t)$.
This is a bijective presentation of the same formulas, not an axiom deletion.
In particular, $k$ is explicit even when some of its bound variables are
absent from both endpoints.

With de Bruijn indices, a term variable is either a negative free level,
zero, or a positive bound index. Opening an external binder of
$(k,s,t)$ at offset $d$ opens its endpoints at depth $2k+d$.
To eliminate the outer $D$, replace $(k,s,t)$ by the matrix
$(k-1,s,t)$ and open its witness at depth $2(k-1)+1$.
The remaining universal coordinate opens at depth $2(k-1)$.
These formulas account for all binder arithmetic in the checker.

\begin{lemma}[Ranked representation]
Let $\rho(D^k(s=t))=(k,\rho_T(s),\rho_T(t))$, where $\rho_T$
erases the term constructor tags and stores successor as a singleton and
iteration as a triple. Then $\rho$ and its evident inverse are mutual
inverses on well-scoped syntax. They commute with opening, including the
contraction of an iterator whose counter becomes literal zero.
\end{lemma}
\begin{proof}
The term assertion follows by structural induction. The only binder in
a numerical term is the last argument of the ternary iterator, so its
depth increases by one. The zero-contraction test is preserved because
literal zero is represented by the same integer. A formula has exactly
$k$ paired binders, contributing depth $2k$. Consequently opening its
endpoints at depth $2k+d$ is exactly the recursive tagged opening through
those binders. The inverse inserts the unique $k$ paired binders and the
term tags determined by arity.
\end{proof}

The certificate translation preserves every rule tag, clause index,
premise list, and proof-tree edge. It changes only the representation of
terms and formulas. Induction on the proof tree therefore gives exact
preservation and reflection of certificate acceptance. This representation
theorem is logically independent of the proof that the four-rule Horn
calculus presents iPRA.

\paragraph{Scope from context length.}
For closed certificates the context starts with one fixed reflexivity
clause. Each introduction or elimination branch adds one clause and one
fresh numerical parameter; no other rule changes either count.
Therefore the number of available free levels is always
$\lvert\Gamma\rvert-1$. The scoped implementation computes this count
instead of passing it through every recursive call. This applies only to
the closed-certificate entry point with its fixed initial context.
It is an implementation simplification, not a change of theory.

\paragraph{Measured comparison.}
After removing only the initial documentation string and applying
CPython 3.12.14's \texttt{ast.unparse} with a final newline, the tagged
four-rule source has 2,888 bytes, 823 significant tokens, and 893 AST nodes.
Ranked equations with arity-coded terms have 2,712 bytes, 789 tokens,
and 857 AST nodes. Deriving scope from context length gives 2,692 bytes,
781 tokens, and 853 AST nodes. Deterministic gzip at compression level
nine has sizes 1,006, 963, and 964 bytes, respectively.
These are matched comparisons with a straightforward checker. They are
not lower than the repository's much more compact 788-byte validated
codec or 516-byte compact profile.

\paragraph{Verification.}
Both ranked implementations accept 86 translated positive certificates
and reject 38 additional malformed certificates. Each run was repeated
with normal Python and with \texttt{-O}. The tests also check 2,000
generated instances of opening/representation coherence and exact
round trips of the 86 certificates. Finite tests support the implementation
review; they are not a proof of arithmetic soundness or bi-interpretability.

\paragraph{Rejected implementation experiment.}
Omitting the public rank and inferring it from endpoint support requires
canonicalizing vacuous outer paired binders and retaining explicit ranks
in internal clause bodies. A support traversal can double as a term
validator, but extra rank checks in Step and local-fact canonicalization
increase the examined checker to 3,069 bytes, 854 tokens, and 957 AST
nodes (1,061 gzip bytes). It is not the selected implementation, and its
full mathematical normalization argument is not claimed here.
